\chapter{Objects or chemicals in the eye}
\index{Objects or chemicals in the eye}

\section*{Definition and Overview}
Ocular trauma involving foreign bodies or chemical exposures represents a critical ophthalmic emergency that can range from minor discomfort to irreversible blindness. Ocular foreign bodies can be classified as superficial (located on the bulbar or palpebral conjunctiva, or embedded in the cornea) or intraocular (penetrating the globe). Chemical injuries represent another form of severe ocular trauma and are broadly divided into acid and alkali burns. The clinical course and ultimate visual outcome are heavily dependent on the type of exposure, the speed of initial decontamination (particularly for chemical injuries), and the precision of subsequent medical intervention.

\section*{Epidemiology}
Ocular foreign bodies and chemical exposures are highly prevalent, accounting for a significant proportion of emergency department visits. They occur most frequently in young men, commonly in occupational settings such as construction, manufacturing, welding, and agriculture. Chemical eye injuries occur equally in domestic and industrial environments, with common household culprits including drain cleaners, bleach, plaster, cement, and car battery fluid. Accidental chemical splashes also represent a common pediatric emergency, often involving liquid laundry detergent capsules or gardening products.

\section*{Aetiology and Pathophysiology}
The biological mechanisms of ocular damage vary depending on the nature of the offending agent:
\begin{itemize}
    \item \textbf{Alkali chemical burns:} Alkalis (e.g., ammonia, sodium hydroxide, lime) are lipophilic and cause rapid liquefactive necrosis. They saponify cell membrane lipids, allowing the chemical to penetrate deep into the anterior chamber, trabecular meshwork, iris, and ciliary body within minutes, leading to severe intraocular damage.
    \item \textbf{Acid chemical burns:} Acids (e.g., hydrochloric acid, sulphuric acid) generally cause coagulative necrosis. They precipitate mucosal proteins, which form a protective physical barrier that limits deeper penetration into the eye. However, strong acids can still cause significant corneal thinning and perforation.
    \item \textbf{Superficial foreign bodies:} Particles such as dust, wood shavings, glass, or metallic shards (often from grinding or drilling) become lodged in the corneal epithelium or palpebral conjunctiva. This triggers localized inflammation, corneal abrasion, and risks secondary bacterial invasion.
    \item \textbf{Intraocular foreign bodies (IOFBs):} High-velocity projectiles (typically metal-on-metal strikes) penetrate the sclera or cornea, entering the vitreous chamber. This carries a high risk of endophthalmitis, retinal detachment, and toxic reactions (e.g., siderosis from iron, chalcosis from copper).
\end{itemize}

\section*{Clinical Features}
Ocular trauma presents with acute, localized symptoms that reflect the severity of the injury:
\begin{itemize}
    \item Sudden, intense ocular pain, burning, or a gritty foreign body sensation
    \item Blepharospasm (involuntary, tight closure of the eyelids)
    \item Epiphora (excessive tearing) and photophobia (sensitivity to light)
    \item Conjunctival injection (redness) and chemosis (swelling of the conjunctiva)
    \item Decreased visual acuity or blurred vision
    \item Corneal haze, opacification, or a chalky white appearance (limbal ischaemia) in severe chemical burns
    \item Subconjunctival haemorrhage or a visible foreign body
    \item Distortion of the pupil or a shallow anterior chamber (suggestive of globe penetration)
\end{itemize}

\section*{Differential Diagnosis}
Ocular pain and redness from foreign bodies or chemicals must be differentiated from:
\begin{itemize}
    \item \textbf{Corneal abrasion:} A scratch on the corneal surface where the foreign body is no longer present.
    \item \textbf{Infectious keratitis:} Bacterial, viral, or fungal corneal ulcers presenting with pain, redness, and purulent discharge.
    \item \textbf{Acute angle-closure glaucoma:} Sudden rise in intraocular pressure causing severe eye pain, headache, nausea, and a mid-dilated, non-reactive pupil.
    \item \textbf{Uveitis:} Intraocular inflammation causing deep, aching pain, photophobia, and ciliary flush.
    \item \textbf{Photokeratitis:} Corneal burns caused by exposure to ultraviolet light (e.g., welder's flash or snow blindness), which presents bilaterally.
\end{itemize}

\section*{When to Seek Medical Care}
All chemical exposures and suspected penetrating eye injuries require immediate, emergency management. For chemical exposures, eye irrigation must begin immediately at the scene before transport.

\textbf{Contact your local emergency services for:}
\begin{itemize}
    \item Any chemical splash to the eye (begin flushing immediately with water or saline)
    \item Suspected penetrating eye injury (e.g., a visible object protruding from the eye, or injury from a high-velocity projectile)
    \item Sudden, severe reduction or total loss of vision
    \item A pupil that appears distorted, tear-shaped, or is leaking fluid
    \item Severe pain accompanied by vomiting or deep headache
\end{itemize}

\section*{Diagnosis and Investigations}
A rapid, systematic assessment is vital to guide treatment while preventing secondary damage:
\begin{itemize}
    \item \textbf{Immediate pH testing (for chemical injuries):} Litmus paper is used to check the pH of the conjunctival fornix. Irrigation must proceed immediately without waiting for other assessments.
    \item \textbf{Visual acuity testing:} The first diagnostic step in non-chemical injuries. Visual acuity must be measured in both eyes individually using a Snellen chart.
    \item \textbf{Slit-lamp biomicroscopy:} Comprehensive examination of the eyelids, conjunctiva, cornea, anterior chamber, and iris to locate foreign bodies and evaluate corneal damage.
    \item \textbf{Eyelid eversion:} The upper and lower eyelids must be everted to inspect the palpebral conjunctiva and fornices for trapped particulate matter (e.g., concrete dust or seeds).
    \item \textbf{Fluorescein staining:} Instillation of fluorescein dye to highlight epithelial defects, corneal abrasions, and to perform Seidel's test (checking for aqueous humour leakage, which confirms globe penetration).
    \item \textbf{Imaging:} Non-contrast orbital CT is the modality of choice for locating suspected intraocular foreign bodies. MRI is strictly contraindicated if metallic foreign bodies are suspected.
\end{itemize}

\section*{Management}
Management priorities focus on immediate decontamination, removal of foreign bodies, and protection of the globe:
\begin{itemize}
    \item \textbf{Copious irrigation (for chemical burns):} Irrigate immediately with sterile saline, Hartmann's solution, or clean tap water. Continue irrigation for at least 15 to 30 minutes (using a Morgan lens if available) until the pH of the eye is neutralised (7.0 to 7.4), checked 5 to 10 minutes after stopping irrigation.
    \item \textbf{Removal of superficial foreign bodies:} Under topical anaesthesia (e.g., 0.5\% tetracaine), superficial particles can be swept away using a sterile cotton-tipped applicator, or lifted using a sterile needle or foreign body spud under slit-lamp guidance. Residual rust rings from metallic objects must be drilled out using an ophthalmic burr.
    \item \textbf{Penetrating injury stabilisation:} Do not attempt to remove a penetrating object. Place a rigid eye shield (which rests on the bony orbit, not the eye) to protect the globe, keep the patient fasting, and refer immediately to an ophthalmologist for surgical repair.
    \item \textbf{Pharmacological therapy:} Topically administered broad-spectrum antibiotics (e.g., chloramphenicol or ciprofloxacin) to prevent infection; cycloplegics (e.g., cyclopentolate) to relieve ciliary spasm and pain; and preservative-free lubricants to promote epithelial healing.
\end{itemize}

\section*{Complications}
Delayed or inadequate management of ocular trauma can result in serious long-term sequelae:
\begin{itemize}
    \item Corneal scarring, neovascularisation, and persistent epithelial defects
    \item Symblepharon (adhesion between the bulbar and palpebral conjunctiva)
    \item Secondary glaucoma (due to trabecular meshwork damage or inflammatory blockage)
    \item Cataract formation (due to chemical penetration or direct lens trauma)
    \item Endophthalmitis (severe intraocular infection)
    \item Corneal perforation and phthisis bulbi (shrinkage and wasting of the eyeball)
\end{itemize}

\section*{Prognosis and Outcomes}
The prognosis for superficial foreign bodies is excellent, with most corneal abrasions healing completely within 24 to 72 hours without permanent visual impairment. For chemical injuries, the prognosis depends on the severity of limbal stem cell ischaemia; mild acid burns usually resolve completely, whereas severe alkali burns frequently result in permanent corneal scarring, conjunctivalisation of the cornea, and blindness. High-velocity intraocular foreign bodies carry a guarded prognosis, depending on the path of entry and the presence of retinal detachment or infection.

\section*{Prevention and Self-Care}
\begin{itemize}
    \item Always wear certified safety goggles, glasses, or face shields when grinding, drilling, sawing, hammering, or using power tools.
    \item Use appropriate eye protection when handling household cleaners, industrial chemicals, pool chlorine, or garden sprays.
    \item Keep domestic chemicals, adhesives, and liquid detergent pods locked in cabinets out of reach of children.
    \item In the event of a chemical splash, flush the eye immediately under a tap or eyewash station for at least 15 minutes before seeking medical attention.
    \item Never rub the eye when a foreign body is suspected, as this can embed the particle deeper or cause extensive corneal scratches.
\end{itemize}

\section*{Sources and Further Reading}
The following organisations provide further information on ocular first aid, chemical safety, and eye care:
\begin{itemize}
    \item Royal Australian and New Zealand College of Ophthalmologists. \textit{Eye Safety and Injury Prevention.} https://ranzco.edu/
    \item Healthdirect Australia. \textit{Eye injuries.} https://www.healthdirect.gov.au/
    \item Mayo Clinic. \textit{Chemical splashes in the eye: First aid.} https://www.mayoclinic.org/
    \item NHS. \textit{What to do if you get something in your eye.} https://www.nhs.uk/
    \item MSD Manual (Professional Version). \textit{Chemical Corneal Burns and Conjunctival Foreign Bodies.} https://www.msdmanuals.com/
\end{itemize}
